% Results and Discussion chapter for:
% A Hierarchical Time-Aware Transformer for Flow-Level Network Intrusion Detection
%
% Required preamble packages:
%   \usepackage{graphicx}
%   \usepackage{booktabs}
% If the thesis is compiled from another directory, redefine this path in the
% preamble, for example: \renewcommand{\resultsfigurepath}{figures}
\providecommand{\resultsfigurepath}{doc/figures}

\chapter{Results and Discussion}
\label{sec:results}

This chapter reports and interprets the empirical results obtained under the
protocol defined in Chapter~\ref{sec:experimental_methodology}. Results and
discussion are integrated within each experiment block so that an observation,
its operational meaning, and its limitations are presented together. The
chapter therefore does not repeat the preprocessing pipeline, model equations,
training procedure, threshold-selection rule, or metric definitions already
given in Chapters~\ref{sec:dataset_preprocessing},
\ref{sec:proposed_methodology}, and
\ref{sec:experimental_methodology}.

Unless stated otherwise, F1, precision, and recall refer to their
macro-averaged values. Class-1 measures are reported separately because label
1 is the alert-derived positive class. The analysis prioritizes macro-F1,
class-1 F1, PR-AUC, and false-positive rate (FPR) over accuracy alone. This is
important because the positive class represents only $3{,}520/84{,}380=4.17\%$
of the principal test manifest; under such imbalance, accuracy and ROC-AUC may
remain high even when positive retrieval is materially weaker
\cite{SaitoRehmsmeier2015}. All confusion matrices are ordered as
TN/FP/FN/TP, and all threshold-dependent metrics use the threshold selected on
the validation set.

\section{Comparison Scope}
\label{subsec:results_comparison_boundaries}

The available experiments use two principal test manifests. Their boundaries must be
preserved because equal sample counts do not imply identical examples or class
composition. Table~\ref{tab:results_manifest_scope} states which comparisons
are controlled.

\begin{table}[htbp]
    \centering
    \caption{Test-manifest boundaries used in this chapter.}
    \label{tab:results_manifest_scope}
    \begin{tabular}{@{}p{5.1cm}rrrp{4.9cm}@{}}
        \toprule
        Experiment block & $N$ & $N_0$ & $N_1$ & Permitted comparison \\
        \midrule
        Stage~1 sequence models, flow-statistics MLP, encodings, sequence
        lengths, packet-selection strategies, and controlled Schemes A--C
        & 35,410 & 34,003 & 1,407 & Direct comparison within this block \\
        Full-data Stage~1 Scheme C reference and all principal Stage~2 experiments
        & 84,380 & 80,860 & 3,520 & Direct comparison within this block \\
        \bottomrule
    \end{tabular}
\end{table}

The flow-statistics MLP now uses the same 34,003/1,407 class composition as the
sequence models and is included as a direct Stage~1 baseline. Two Scheme C
evaluations are retained for different purposes: the 35,410-flow run provides
the controlled A--C fusion comparison, whereas the separate 84,380-flow run
supplies the frozen embeddings and independent Stage~1 baseline used in the
Stage~2 analysis. Results from these two manifests are never mixed in one
effect-size calculation. The capture is also
organization-specific, so numerical comparisons with studies using public
datasets, different attack taxonomies, or different labeling rules would not
be controlled. These restrictions follow established guidance on avoiding
dataset and evaluation artefacts in intrusion-detection research
\cite{SommerPaxson2010,Arp2022DosDonts}.

\section{Stage~1: Intra-Flow Representation}
\label{subsec:results_stage1}

\subsection{Architecture and Temporal-Encoding Effects}
\label{subsubsec:results_stage1_models}

Figure~\ref{fig:results_stage1_models} compares the Stage~1 model families and
encoding ablations on the common 35,410-flow test set. A figure is used here
instead of a wide predictive-performance table because the main question is
the relative pattern across macro-F1, class-1 F1, and PR-AUC. The exact
computational measurements are retained separately in
Table~\ref{tab:results_stage1_errors_efficiency}.

\begin{figure}[htbp]
    \centering
    \includegraphics[width=\textwidth]{\resultsfigurepath/stage1_model_comparison.pdf}
    \caption{Stage~1 architecture and encoding comparison on the common
    35,410-flow test manifest. The horizontal axis begins at 0.52 to make the
    differences visible; absolute values are printed at the end of each bar.}
    \label{fig:results_stage1_models}
\end{figure}

The combined positional-and-cumulative-time encoder produces the strongest
overall result: macro-precision $0.8896$, macro-recall $0.8869$, macro-F1
$0.8882$, class-1 F1 $0.7853$, ROC-AUC $0.9860$, and PR-AUC $0.8564$. Its
confusion matrix is $33{,}707/296/306/1{,}101$, corresponding to an FPR of
$0.8705\%$. Relative to LSTM+time, the strongest non-Transformer baseline by
macro-F1, it gains $0.0205$ macro-F1, $0.0398$ class-1 F1, and $0.0487$
PR-AUC. It also recovers 95 additional positive flows, while adding only 10
false positives.

The encoding ablation shows that ordinal position and elapsed time are
complementary rather than interchangeable. Compared with no encoding, the
combined representation increases macro-F1 by $0.0375$ and PR-AUC by
$0.0561$, while reducing FP from 420 to 296 and FN from 394 to 306. Time-only
encoding is stronger than position-only encoding, particularly in macro-recall
($0.8776$ versus $0.8642$), but combining them gives a further $0.0064$
macro-F1 improvement over time only. Standard positional encoding represents
token order rather than physical elapsed time \cite{Vaswani2017}; this result
is therefore consistent with traffic-specific evidence that timing carries
additional information \cite{Han2023GTID}.

The lowest FPR in this block belongs to the GRU ($0.8293\%$), but it also has
437 FN, the largest missed-positive count among these models. The result
illustrates why FPR cannot be interpreted in isolation: fewer alerts may simply
reflect a more conservative detector. The proposed Stage~1 model instead
provides the best macro-F1 and class-1 F1, indicating a better balance between
false alerts and missed attacks.

The flow-statistics MLP establishes the value of packet-sequence modelling on
the same test composition. Metrics recomputed from its confusion matrix are
macro-precision $0.7501$, macro-recall $0.8094$, macro-F1 $0.7762$, class-1 F1
$0.5723$, and FPR $2.5086\%$; the run reports ROC-AUC $0.9552$ and PR-AUC
$0.5604$. The position+time Transformer raises macro-F1 by $0.1121$, class-1
F1 by $0.2130$, ROC-AUC by $0.0308$, and PR-AUC by $0.2960$. Its confusion
matrix also contains 557 fewer FP and 195 fewer FN. Hence the advantage is not
limited to threshold-free ranking: the packet-sequence model simultaneously
reduces false-alert burden and missed positives at its selected operating
point.

Figure~\ref{fig:results_stage1_errors} makes the operating-point trade-offs
explicit. Because every model has the same 34,003 negative examples, the FP
bars are directly proportional to FPR. GRU has the lowest FPR, but its 437 FN
show that the result is achieved with weak positive coverage. By contrast, the
position+time model adds only 14 FP relative to GRU while removing 131 FN.

\begin{figure}[htbp]
    \centering
    \includegraphics[width=\textwidth]{\resultsfigurepath/stage1_error_comparison.pdf}
    \caption{Stage~1 false-positive and false-negative counts on the common
    35,410-flow test manifest. FPR is printed beside each FP count.}
    \label{fig:results_stage1_errors}
\end{figure}

Explicit time information also benefits the recurrent baseline. LSTM+time
raises macro-F1 from $0.8582$ to $0.8677$, class-1 F1 from $0.7273$ to
$0.7455$, and PR-AUC from $0.7976$ to $0.8077$, while removing 29 FP and 22
FN. The CNN result is more nuanced: CNN1D+time improves macro-F1 and class-1 F1
and removes 116 FP, but its PR-AUC falls from $0.7975$ to $0.7821$ and FN rises
by 30. Thus, an improved selected operating point does not necessarily imply
better probability ranking across all thresholds.

\begin{table}[htbp]
    \centering
    \caption{Stage~1 parameter counts and single-run training times measured on
    a Google Colab T4 GPU.}
    \label{tab:results_stage1_errors_efficiency}
    \begin{tabular}{@{}lrrr@{}}
        \toprule
        Model & Trainable parameters & Training time (s) & Training time (min) \\
        \midrule
        Position+time Transformer & 344,963 & 4,337.6 & 72.29 \\
        Time-only Transformer & 344,963 & 2,250.5 & 37.51 \\
        Position-only Transformer & 344,963 & 1,994.7 & 33.25 \\
        LSTM+time & 217,346 & 3,556.3 & 59.27 \\
        CNN1D+time & 70,274 & 1,162.4 & 19.37 \\
        GRU & 162,690 & 3,684.4 & 61.41 \\
        LSTM & 216,834 & 3,381.2 & 56.35 \\
        CNN1D & 68,354 & 1,591.0 & 26.52 \\
        BiLSTM & 564,738 & 2,625.7 & 43.76 \\
        No-encoding Transformer & 344,963 & 2,213.9 & 36.90 \\
        \bottomrule
    \end{tabular}
\end{table}

The MLP is omitted from Table~\ref{tab:results_stage1_errors_efficiency}
because its parameter count and training time were not included in the
supplied run summary.

CNN1D+time is the fastest Stage~1 model at 1,162.4 seconds, whereas the
position+time Transformer takes 4,337.6 seconds. The proposed encoder is not
the largest model, because the BiLSTM contains 564,738 parameters, yet it is
the most expensive recorded run. Because all Stage~1 runs in this block were
executed on a Colab T4 GPU, the hardware class is controlled. The values still
represent total single-run training cost rather than pure per-epoch throughput:
model-specific early stopping, data loading, and Colab runtime variation can
affect wall-clock time. Parameter count alone therefore does not explain the
observed training duration.

\subsection{Packet and Flow-Statistics Integration}
\label{subsubsec:results_stage1_fusion}

All three schemes in Figure~\ref{fig:results_stage1_fusion} use the same
35,410-flow test manifest. Scheme A tiles the flow-statistics vector into every
packet token, Scheme B uses packet features only, and Scheme C keeps packet and
flow representations separate before bounded token-FiLM fusion.

\begin{figure}[htbp]
    \centering
    \includegraphics[width=\textwidth]{\resultsfigurepath/stage1_fusion_comparison.pdf}
    \caption{Controlled Stage~1 integration comparison on the common
    35,410-flow manifest. Left: macro metrics and threshold-independent ranking
    metrics. Right: false-positive and false-negative counts. The score axis is
    truncated and explicitly labelled.}
    \label{fig:results_stage1_fusion}
\end{figure}

Within the controlled A--B comparison, Scheme B increases macro-F1 by
$0.0130$, macro-recall by $0.0269$, ROC-AUC by $0.0059$, and PR-AUC by
$0.0238$. It detects 77 additional positives, reducing FN from 383 to 306,
while producing 32 more FP. Scheme A has slightly higher macro-precision
($0.8919$ versus $0.8896$), but its weaker recall and ranking performance show
that repeating the same flow-statistics vector at every packet position does
not replace sequence-specific evidence. This supports keeping packet and
statistical representations in distinct branches before controlled fusion.

Scheme C obtains the strongest controlled result: macro-precision $0.8990$,
macro-recall $0.8920$, macro-F1 $0.8955$, class-1 F1 $0.7991$, ROC-AUC
$0.9890$, and PR-AUC $0.8763$. Its confusion matrix is
$33{,}736/267/293/1{,}114$, giving FPR $0.7852\%$ and FNR $20.8244\%$.
Relative to Scheme B, C raises macro-F1 by $0.0072$, class-1 F1 by $0.0138$,
and PR-AUC by $0.0199$, while removing 29 FP and 13 FN. FPR consequently falls
from $0.8705\%$ to $0.7852\%$. This is the cleanest evidence that the separate
flow-statistics branch contributes information beyond the packet-only encoder,
because both error types improve on the same test examples.

The comparison with Scheme A is operationally even clearer. Scheme C misses 90
fewer positive flows while producing only three additional false alerts. Its
FPR is therefore nearly unchanged ($0.7852\%$ versus $0.7764\%$), while
macro-F1 increases by $0.0202$ and PR-AUC by $0.0436$. Tiling an aggregate
vector into every token appears less effective than preserving separate
packet- and flow-level pathways and allowing a bounded modulation mechanism to
combine them after sequence encoding.

The validation-selected thresholds differ substantially: 0.715 for A, 0.455
for B, and 0.683 for C. These values should not be interpreted as directly
comparable confidence scales because the three models need not be equally
calibrated. The ordering is nevertheless supported by ROC-AUC and PR-AUC,
which do not depend on a single operating threshold. Scheme B has the lowest
test loss ($0.0576$ versus $0.0684$ for C), but the principal selection
criterion is macro-F1; loss also reflects probability calibration and the
training objective rather than the final error balance alone.

For the later hierarchical experiment, a separate full-data Scheme C run is
evaluated on 84,380 flows and obtains macro-F1 $0.8954$, PR-AUC $0.8763$, FPR
$0.8249\%$, and confusion matrix $80{,}193/667/732/2{,}788$. That run is used
only as the Stage~1 baseline in Table~\ref{tab:results_hierarchical_gain}; its
counts are not substituted into the controlled A--C comparison above.

\subsection{Sequence Length and Packet Selection}
\label{subsubsec:results_sequence_length}

Figure~\ref{fig:results_stage1_sensitivity} combines the two Stage~1
sensitivity studies. The line plot reveals whether additional packet history
produces monotonic gains, while the grouped bars expose the precision--recall
trade-off among packet-selection strategies without another wide table.

\begin{figure}[htbp]
    \centering
    \includegraphics[width=\textwidth]{\resultsfigurepath/stage1_sensitivity.pdf}
    \caption{Stage~1 sensitivity analysis. Left: maximum packet-sequence length
    under head selection. Right: packet-selection strategy at $L=128$. Dashed
    lines identify the selected sequence length. Truncated axes are explicitly
    labelled.}
    \label{fig:results_stage1_sensitivity}
\end{figure}

Performance is non-monotonic in sequence length. Short inputs ($L=8$ and
$L=16$) produce the weakest PR-AUC values, $0.8369$ and $0.8455$, suggesting
that the initial packets alone do not always contain sufficient ranking
evidence. At $L=32$, PR-AUC reaches its maximum of $0.8571$, but macro-F1 is
only $0.8817$. At $L=64$, macro-recall is highest ($0.8907$) and FN is lowest
(294), but FPR rises to $0.9676\%$. The selected $L=128$ configuration gives
the best macro-F1 ($0.8882$), class-1 F1 ($0.7853$), ROC-AUC ($0.9860$), and
the lowest FPR in this length sweep ($0.8705\%$). Increasing the limit to 256
reduces macro-F1 by $0.0062$ and PR-AUC by $0.0071$, while adding 46 FP.
Consequently, 128 packets represent a balanced evidence budget rather than an
arbitrary maximum.

Head and head-tail selection are practically close. Head improves macro-F1 by
$0.0016$ and produces eight fewer FN, whereas head-tail removes one FP and has
a $0.0002$ PR-AUC advantage. Such differences are too small for a strong
single-seed claim. Random selection obtains the highest macro-precision
($0.8959$) and lowest FPR ($0.7646\%$), but macro-recall falls to $0.8707$ and
FN increases to 353. Head selection is retained because it has the strongest
macro-F1, class-1 F1, and positive coverage while remaining deterministic and
compatible with causal deployment.

\subsection{Stage~1 Answer to RQ1}
\label{subsubsec:results_rq1}

The controlled Stage~1 results support RQ1. Packet order and cumulative time
provide complementary information, and their combination outperforms the
recurrent and convolutional baselines by macro-F1 and PR-AUC. The result is not
explained solely by parameter count because the larger BiLSTM performs worse.
The sensitivity analysis also qualifies the finding: the advantage depends on
a moderate packet budget and stable selection rule; simply providing more
packets does not guarantee better detection.

\section{Stage~2: Inter-Flow Context}
\label{subsec:results_stage2}

\subsection{Gain over Independent Flow Classification}
\label{subsubsec:results_hierarchical_gain}

Table~\ref{tab:results_hierarchical_gain} gives the principal hierarchical
comparison. Both rows use exactly the same 84,380-flow test manifest. The first
row classifies the Scheme C embedding independently, while the second adds
causal source-host context to the frozen Stage~1 representation.

\begin{table}[htbp]
    \centering
    \caption{Independent Stage~1 classification and the complete hierarchical
    system on the common full-data manifest.}
    \label{tab:results_hierarchical_gain}
    \resizebox{\textwidth}{!}{%
    \begin{tabular}{@{}lrrrrrrrl@{}}
        \toprule
        Model & $P_{\mathrm{macro}}$ & $R_{\mathrm{macro}}$ &
        $F_{1,\mathrm{macro}}$ & $F_{1,1}$ & ROC-AUC & PR-AUC & FPR (\%) &
        TN/FP/FN/TP \\
        \midrule
        Stage~1 Scheme C & 0.8990 & 0.8919 & 0.8954 & 0.7994 & 0.9890 &
        0.8763 & 0.8249 & 80,193/667/732/2,788 \\
        Stage~2 source-host & \textbf{0.9328} & \textbf{0.9224} &
        \textbf{0.9275} & \textbf{0.8610} & \textbf{0.9938} &
        \textbf{0.9345} & \textbf{0.5429} & 80,421/439/527/2,993 \\
        \bottomrule
    \end{tabular}%
    }
\end{table}

Adding source-host context raises macro-F1 by $0.0321$, macro-precision by
$0.0338$, macro-recall by $0.0305$, class-1 F1 by $0.0616$, and PR-AUC by
$0.0582$. It simultaneously removes 228 FP and 205 FN. FPR falls from
$0.8249\%$ to $0.5429\%$, a relative reduction of approximately $34.2\%$.
Because both error types improve on the same flows, the gain is not merely a
threshold-induced exchange between precision and recall. The smaller ROC-AUC
increase ($0.0048$) is expected in this imbalanced setting; PR-AUC exposes the
improvement in positive retrieval more clearly.

\subsection{Which Relation Defines Useful Context?}
\label{subsubsec:results_context_relation}

Figure~\ref{fig:results_stage2_context} combines ranking performance with the
operational recall--FPR trade-off. The left panel compares macro-F1 and PR-AUC;
the right panel places each context relation according to class-1 recall and
false-positive rate.

\begin{figure}[htbp]
    \centering
    \includegraphics[width=\textwidth]{\resultsfigurepath/stage2_context_analysis.pdf}
    \caption{Stage~2 context-relation ablation at $W=128$. Left: macro-F1 and
    PR-AUC. Right: class-1 recall versus FPR, where the desirable region is the
    upper-left corner.}
    \label{fig:results_stage2_context}
\end{figure}

Source-host context provides the strongest overall balance, with macro-F1
$0.9275$, PR-AUC $0.9345$, class-1 F1 $0.8610$, ROC-AUC $0.9938$, and the
lowest FPR ($0.5429\%$). It exceeds destination-host and time-only context by
$0.0079$ and $0.0074$ macro-F1, respectively, while reducing both FP and FN in
both comparisons. This suggests that recent behavior associated with the
initiating host is more discriminative than global recency or destination
identity alone. The result is consistent with work showing that relational or
window-level traffic structure can reveal behavior omitted by independent-flow
classifiers \cite{Deng2023FlowTopology,MartinReizabal2026FlowWindow}.

Endpoint context occupies a different operating regime. It has the highest
macro-recall ($0.9387$) and class-1 recall ($0.8861$), missing only 401
positives, but it also produces 702 FP and the highest FPR ($0.8682\%$).
Source-host context removes 263 of those false alerts and increases macro-F1 by
$0.0061$, at the cost of 126 additional FN. Endpoint context may therefore be
preferable when missed attacks are far more costly than analyst workload;
source-host context is the better default under the thesis's macro-F1
criterion. Time-only context has the second-lowest FPR but the weakest PR-AUC
and the largest FN count, showing that temporal proximity without an entity
relation is insufficient.

\subsection{How Much Context Is Useful?}
\label{subsubsec:results_context_window}

\begin{figure}[htbp]
    \centering
    \includegraphics[width=\textwidth]{\resultsfigurepath/stage2_window_sensitivity.pdf}
    \caption{Source-host context-window sensitivity. Left: macro-F1 and
    PR-AUC. Right: false-positive and false-negative counts. The dashed line
    marks the selected window $W=128$.}
    \label{fig:results_stage2_window}
\end{figure}

As shown in Figure~\ref{fig:results_stage2_window}, the response to context
length is non-monotonic. $W=32$ gives the highest macro-recall ($0.9310$) and
only 452 FN, but also generates 775 FP and the highest FPR ($0.9584\%$). At
$W=128$, macro-F1 rises to $0.9275$ and PR-AUC to $0.9345$, while FP falls to
439. Relative to $W=256$, the selected window removes 121 FP and 45 FN and
improves macro-F1 by $0.0116$. Thus, more history is not automatically more
informative: older or weakly related source-host flows can dilute the evidence,
whereas a window of 128 retains sufficient recent context without excessive
noise.

\subsection{Stage~2 Architecture Comparison}
\label{subsubsec:results_stage2_architecture_comparison}

The architecture comparison holds the Stage~1 embeddings, source-host
relation, $W=128$, online causal policy, training objective, optimizer,
sampling rule, checkpoint criterion, and validation threshold search fixed.
The model-specific encoder and its dimensions are changed. The six models are
presented graphically in Figure~\ref{fig:results_stage2_architectures}; the
second panel displays FP and FN directly so that a higher F1 score cannot hide
an unfavorable error trade-off.

\begin{figure}[htbp]
    \centering
    \includegraphics[width=\textwidth]{\resultsfigurepath/stage2_architecture_comparison.pdf}
    \caption{Stage~2 architecture comparison on the common 84,380-flow test
    manifest. Left: macro-F1, class-1 F1, and PR-AUC. Right: threshold-dependent
    false-positive and false-negative counts. Truncated score axes are labelled.}
    \label{fig:results_stage2_architectures}
\end{figure}

The proposed target-query gated model is best in macro-precision ($0.9328$),
macro-F1 ($0.9275$), class-1 F1 ($0.8610$), ROC-AUC ($0.9938$), PR-AUC
($0.9345$), test loss ($0.019559$), and FPR ($0.5429\%$). CNN+LSTM is the
strongest baseline by macro-F1 ($0.9097$), GRU by PR-AUC ($0.9141$), and LSTM
by macro-recall ($0.9243$). The proposed model improves over these respective
best baseline values by $0.0179$ macro-F1 and $0.0204$ PR-AUC, while its
macro-recall is only $0.0019$ below LSTM.

The confusion matrices clarify this balance. Compared with the vanilla
Transformer, the proposed model removes 177 FP and 117 FN. Compared with GRU,
it removes 127 FP and 116 FN; compared with CNN+LSTM, it removes 75 FP and 150
FN. LSTM detects 28 more positives than the proposed model, but produces 338
additional false alerts. Therefore, the proposed model's F1 gain does not come
from sacrificing positive coverage for precision.

The comparison with vanilla self-attention also supports the target-query
design. Vanilla attention updates the complete context sequence and classifies
the last valid token, whereas the proposed model uses the current flow as the
query, summarizes eligible history, and injects the update through a learned
feature-wise gate. The proposed model gains $0.0213$ macro-F1, $0.0407$
class-1 F1, and $0.0272$ PR-AUC over vanilla attention. This supports, but does
not alone prove, the value of the inductive bias because their parameter counts
differ.

The no-context head provides a current-flow-only lower bound. The proposed
model improves macro-F1 by $0.0331$, class-1 F1 by $0.0631$, and PR-AUC by
$0.0559$, while reducing FPR from $1.0648\%$ to $0.5429\%$. However, the
no-context head has only 514 parameters, compared with 283,522 for the proposed
model. A capacity-matched confirmatory ablation should retain the target-query
architecture while masking all historical tokens before attributing the full
gain causally to context.

\paragraph{PR-AUC versus a precision--recall curve.}
PR-AUC is a scalar area and cannot itself be drawn as a curve. The left panel
of Figure~\ref{fig:results_stage2_architectures} therefore presents PR-AUC as
bars. A genuine precision--recall curve requires every test flow's label and
predicted probability across thresholds. The aggregate PR-AUC values and
confusion matrices are insufficient to reconstruct that curve. The final
experiment archive should retain aligned \texttt{stage2\_predictions\_test.csv}
files for all six models; when these files are available, the following figure
is included automatically.

\IfFileExists{\resultsfigurepath/stage2_pr_curves.pdf}{%
\begin{figure}[htbp]
    \centering
    \includegraphics[width=0.78\textwidth]{\resultsfigurepath/stage2_pr_curves.pdf}
    \caption{Precision--recall curves for the Stage~2 architecture comparison.
    Each legend entry reports average precision on the same aligned test flows.}
    \label{fig:results_stage2_pr_curves}
\end{figure}
}{}

\subsection{Efficiency and Operational Trade-Offs}
\label{subsec:results_statistical_efficiency}

Figure~\ref{fig:results_stage2_efficiency} relates parameter count to measured
model-evaluation latency and annotates macro-F1. A logarithmic parameter axis
is used because model sizes range from 514 to more than 1.2 million parameters.

\begin{figure}[htbp]
    \centering
    \includegraphics[width=0.88\textwidth]{\resultsfigurepath/stage2_efficiency.pdf}
    \caption{Stage~2 performance--efficiency trade-off. Latency covers model
    evaluation after Stage~1 embeddings and context indices have been built; it
    is not end-to-end PCAP-to-alert latency.}
    \label{fig:results_stage2_efficiency}
\end{figure}

The no-context head is the smallest and fastest model, with 514 parameters and
0.1812~ms per flow, but it also has the weakest macro-F1 and PR-AUC. The
proposed model uses 283,522 parameters and 0.1965~ms per flow, corresponding to
approximately 5,089 flows per second. Its latency is effectively identical to
the vanilla Transformer's 0.1967~ms; the 0.0002~ms difference is too small for
a practical speed claim. This profile is consistent with computing one
target-to-history query instead of updating every context token.

Relative to GRU, LSTM, and CNN+LSTM, the proposed model uses $59.0\%$,
$69.3\%$, and $76.7\%$ fewer parameters, respectively. Its measured latency is
lower by $42.2\%$, $43.3\%$, and $45.2\%$. These values describe the model
forward pass only. They constitute controlled efficiency comparisons only if
all profiles were collected on the same accelerator and software stack;
otherwise, a hardware-matched T4 or L4 timing run is required. Stage~2 training
time, context-index construction, Stage~1 inference, CSV extraction, and PCAP
processing remain outside the reported latency.

\section{Integrated Findings and Research Questions}
\label{subsec:results_rq_answers}

\paragraph{RQ1: intra-flow temporal representation.}
The results support the claim that packet-sequence modeling benefits from both
order and elapsed time. The combined encoder has the highest macro-F1 and
PR-AUC in the controlled Stage~1 architecture block, and removing either
encoding weakens performance. The selected 128-packet head strategy provides
the strongest overall error balance. The controlled fusion experiment further
shows that bounded token-FiLM integration improves over both tiled flow
statistics and the packet-only branch, reaching macro-F1 $0.8955$ and PR-AUC
$0.8763$ while reducing both FP and FN relative to Scheme B. On the same class
composition, the position+time packet model also substantially outperforms the
flow-statistics MLP. The conclusion is still specific to this organizational
capture and does not by itself establish cross-dataset generalization.

\paragraph{RQ2: inter-flow contextual reasoning.}
The current evidence supports the usefulness of causal cross-flow context.
Relative to independent Stage~1 classification on the same full manifest,
source-host Stage~2 context raises macro-F1 from $0.8954$ to $0.9275$, raises
PR-AUC from $0.8763$ to $0.9345$, and reduces both FP and FN. The target-query
gated architecture also outperforms no-context, vanilla Transformer, LSTM,
GRU, and CNN+LSTM alternatives. Source-host relation and $W=128$ provide the
best balanced operating point, while endpoint context remains a plausible
high-recall alternative.

\section{Reliability and Validity of the Observed Differences}
\label{subsec:results_reliability}

The reported experiments are predominantly single-seed runs using seed 130.
They support ablation and sensitivity observations, but they do not estimate
training variability. Neural-model rankings can vary with initialization and
mini-batch order \cite{ReimersGurevych2017,Bouthillier2021Variance}; small gaps,
such as the $0.0016$ macro-F1 difference between head and head-tail selection,
must therefore be treated as descriptive rather than statistically reliable.

Confidence intervals, paired bootstrap tests, and McNemar tests cannot be
reconstructed from aggregate metrics. They require aligned per-flow
probabilities and predictions for every model. The final confirmatory study
should run three to five seeds, archive each validation-selected threshold,
retain test probabilities, and report mean, standard deviation, and paired
uncertainty intervals. The resolved YAML for the vanilla Transformer and exact
accelerator identity for every efficiency run should also be archived. Until
then, the larger Stage~2 gains are strong empirical observations, but no claim
of statistical significance is made.

\section{Chapter Summary}
\label{subsec:results_summary}

Stage~1 shows that packet position and cumulative time encode complementary
intra-flow evidence. The position+time Transformer reaches macro-F1 $0.8882$
and PR-AUC $0.8564$ on the controlled 35,410-flow manifest, with $L=128$ and
head selection providing the most balanced configuration. On that same
manifest, separate bounded token-FiLM fusion further raises macro-F1 to
$0.8955$ and PR-AUC to $0.8763$, with only 267 FP and 293 FN. On the separate
full 84,380-flow manifest, causal source-host context produces the largest
observed hierarchical gain: macro-F1 increases by 3.21 percentage points,
PR-AUC by 5.82 points, and FPR falls from $0.8249\%$ to $0.5429\%$ while FN
also decreases.

The architecture benchmark further shows that target-query gated attention
outperforms no-context, vanilla Transformer, LSTM, GRU, and CNN+LSTM controls
by macro-F1 and PR-AUC. It achieves nearly the same model-level latency as the
vanilla Transformer and is substantially smaller and faster than the recurrent
alternatives. These findings support the hierarchical design while also
identifying the remaining evidence needed for the final thesis: repeated-seed
statistics, a capacity-matched no-context ablation, aligned prediction files
for PR curves and paired tests, and hardware-matched end-to-end timing.
